\section{Counting and sampling bounded integral flows}\label{sec:flows}

The table theorem also applies to integral single-commodity flows in
arbitrary explicitly listed finite directed networks. Let $G=(V,E)$ be
a directed multigraph, allowing loops, and give each arc finite bounds
$\ell_e,u_e\in\mathbb Z$ and each vertex a balance $\beta_v\in\mathbb Z$,
all encoded in binary. With the convention ``outflow minus inflow equals
balance,'' define
\begin{equation}\label{eq:flow-set}
\begin{aligned}
\mathcal F=\bigl\{f\in\mathbb Z^E:\;&
 \ell_e\le f_e\le u_e\quad(e\in E),\\
 &
 \sum_{e:\,\mathrm{tail}(e)=v}f_e-
 \sum_{e:\,\mathrm{head}(e)=v}f_e=\beta_v\quad(v\in V)\bigr\}.
\end{aligned}
\end{equation}
A loop occurs once in each sum and cancels from the balance. Each
arc-value vector is counted once, not each path or cycle decomposition.
Circulations have $\beta=0$; for distinct $s,t$, prescribed value $k$
from $s$ to $t$ is expressed by balances $k,-k,0$ at $s,t$, and the
other vertices, respectively.

\begin{corollary}[Bounded integral flows]\label{cor:integral-flows}
For the network data above, there are algorithms with these guarantees.
\begin{enumerate}[label=(\roman*)]
\item Given rational $\varepsilon,\delta\in(0,1)$, the counting algorithm
outputs a nonnegative rational $\widehat Z_{\mathrm{flow}}$ satisfying
\[
 \P\bigl((1-\varepsilon)|\mathcal F|\le\widehat Z_{\mathrm{flow}}
             \le(1+\varepsilon)|\mathcal F|\bigr)\ge1-\delta.
\]
If $\mathcal F$ is empty, it outputs zero on every execution.
\item Given rational $\tau\in(0,1)$, the sampling algorithm reports
infeasibility exactly when $\mathcal F$ is empty. Otherwise it outputs
$\widehat f\in\mathcal F$ on every execution, with
\[
 \TV\bigl(\mathcal L(\widehat f),\operatorname{Unif}(\mathcal F)\bigr)
 =\frac12\sum_{f\in\mathcal F}
   \left|\P(\widehat f=f)-\frac1{|\mathcal F|}\right|\le\tau.
\]
\end{enumerate}
Both algorithms use independent unbiased bits. Let $L$ denote the total
binary input length, including the corresponding accuracy parameters.
Fixed polynomials in $L,\varepsilon^{-1},\log(\delta^{-1})$ for counting
and in $L,\tau^{-1}$ for sampling bound the total bit operations on
every execution.
\end{corollary}

\begin{proof}
\emph{The flow-to-table bijection.}
If some $\ell_e>u_e$, or if $\sum_v\beta_v\ne0$, the flow set is empty;
return zero for counting and report infeasibility for sampling. The
second test is necessary because every arc cancels in the sum of the
balance equations. If $V=E=\varnothing$, there is one empty flow, which
can be counted and returned directly. In all other cases, shift the
lower bounds by putting
\[
\begin{aligned}
 x_e&=f_e-\ell_e,\qquad c_e=u_e-\ell_e,\\
 \beta'_v&=\beta_v-\sum_{e:\,\mathrm{tail}(e)=v}\ell_e
                    +\sum_{e:\,\mathrm{head}(e)=v}\ell_e.
\end{aligned}
\]
Then $0\le x_e\le c_e$ and the new balance is $\beta'_v$.

For each original arc $e=(u,v)$ introduce a private new vertex $w_e$
and replace the arc by $u\to w_e\to v$, giving both arcs cap $c_e$ and
giving $w_e$ balance zero. The two arc values must agree at $w_e$, so
this is a bijection on flows, including when $u=v$. The private vertices
make all subdivided arc cells distinct and off-diagonal, also for parallel
arcs. Write $W=V\cup\{w_e:e\in E\}$ and $E'$ for the subdivided vertices
and arcs, and let $c_a$ be the inherited cap of $a\in E'$. Extend
$\beta'$ by zero at the new vertices. For $q=|W|=|V|+|E|>0$, put
\[
\begin{gathered}
 K=1+\sum_{a\in E'}c_a+\sum_{w\in W}|\beta'_w|,\\
 R_w=K+(\beta'_w)^+,\qquad C_w=K+(\beta'_w)^-,
\end{gathered}
\]
where $h^+=\max\{h,0\}$ and $h^-=\max\{-h,0\}$. Index a $q\times q$
table by $W$ in both coordinates, use row margins $R$ and column margins
$C$, and set
\[
 b_{st}=\begin{cases}
 c_a,&a=(s,t)\in E',\\
 K,&s=t,\\
 0,&\text{otherwise}.
 \end{cases}
\]
The cases are disjoint, and the margin totals agree because
$\sum_w\beta'_w=0$.

Given a subdivided flow $z$, put its arc values in the arc cells. At
vertex $w$ the unique possible diagonal entry is
\[
 y_w=R_w-\operatorname{out}_w(z)
     =C_w-\operatorname{in}_w(z).
\]
It is nonnegative since $K$ dominates every incident capacity sum.
If $\beta'_w\ge0$, then $y_w=K-\operatorname{in}_w(z)$; if
$\beta'_w\le0$, then $y_w=K-\operatorname{out}_w(z)$. Thus $y_w\le K$,
so the table is feasible. Conversely, subtracting a table's column
equation from its row equation gives the subdivided balance at every
vertex. The private-vertex balances equate the two copies of each
original arc, and the diagonal entries are forced by the displayed
formula. Adding back $\ell$ gives the unique original flow. This proves
a bijection with $\Omega(R,C,b)$.

The table has $q^2$ explicitly written caps. Its parameters are sums
and differences of polynomially many binary integers, so their bit
lengths and the total table-data length $L_{\mathrm{tab}}$ are polynomial
in the network-data length. Both directions of the bijection use only
polynomially many operations on such integers; no capacity is expanded
in unary. The deterministic feasibility test in Section~\ref{sec:scales}
handles all remaining infeasibility, including edgeless and zero-capacity
cases. If the table is infeasible, return zero for counting and report
infeasibility for sampling. Otherwise Theorem~\ref{thm:main} on the table
gives the counting assertion with the stated bit bound.

\emph{Sampling by conditional counts.}
We use the self-reduction principle connecting approximate counting
and almost-uniform generation~\cite[Theorem~6.3]{JVV1986}, implemented
here by bisecting binary cell intervals. The argument below gives the
required feasibility, total-variation, and bounded-bit guarantees.

After this root feasibility test, write $\Omega(r,c,b)=\Omega(R,C,b)$
for the constructed feasible table set and initialize each cell interval
to $[0,b_{ij}]$. At a reached node, let
$[a_{ij},d_{ij}]\subseteq[0,b_{ij}]$ be the current integer interval for
cell $(i,j)$.
Shifting $X_{ij}$ by $a_{ij}$ is a bijection to a bounded-table instance
with
\[
 b'_{ij}=d_{ij}-a_{ij},\qquad
 r'_i=r_i-\sum_j a_{ij},\qquad c'_j=c_j-\sum_i a_{ij}.
\]
Reject negative residual data before applying the feasibility test.
Process the cells in a fixed order. For a nonsingleton interval $[a,d]$,
put $m=\lfloor(a+d)/2\rfloor$ and use the disjoint children $[a,m]$ and
$[m+1,d]$, continuing until every interval is a singleton. A path has at
most
\[
 D=\sum_{i,j}\left\lceil\log_2(b_{ij}+1)\right\rceil
 \le L_{\mathrm{tab}}
\]
decisions, and every conditioned input has polynomial binary length.
If $D=0$, the already-tested feasible table is unique and is returned
directly. No assumption that attainable cell values form an interval
is used.

At a reached node, test both children for feasibility. Their feasible
sets partition the parent's feasible set, so at least one is feasible.
If exactly one is feasible, take it. If both are feasible, let their positive counts
be $N_0,N_1$ and call Theorem~\ref{thm:main} separately on the two child
instances with relative error $\eta$ and failure probability $\gamma$,
using fresh independent bits. Write the nonnegative outputs as
$\widehat N_0,\widehat N_1$. If their sum is zero, take child zero,
which is known to be feasible. Otherwise put
$\widehat p=\widehat N_0/(\widehat N_0+\widehat N_1)$. Using $t$ further fresh
bits, draw $U$ uniformly from $\{0,\ldots,2^t-1\}$ and take child zero
exactly when $U<\lfloor2^t\widehat p\rfloor$. This implements probability
$p_t=2^{-t}\lfloor2^t\widehat p\rfloor$, with $|p_t-\widehat p|\le2^{-t}$, in bounded
bit time. In particular, a zero estimate is never used as an
infeasibility certificate, and every reached child is feasible.

For exact counts the branch probability is $p=N_0/(N_0+N_1)$, and
these exact conditional probabilities telescope to the uniform law on
feasible singleton leaves. On successful estimates, write
$\widehat N_i=N_i(1+e_i)$ with $|e_i|\le\eta$. Then
\[
 |\widehat p-p|=\frac{p(1-p)|e_0-e_1|}{1+pe_0+(1-p)e_1}
       \le\frac{\eta}{2(1-\eta)}.
\]
At each reached history, the fresh count calls fail with probability at
most $2\gamma$ in total. Averaging over them and the fresh branch bits,
the conditional total-variation error of the implemented branch choice
relative to the exact one is therefore at most
\[
 a=\frac{\eta}{2(1-\eta)}+2\gamma+2^{-t}.
\]
This bound also covers the zero-sum fallback; a one-feasible-child node
has zero error. Couple the two branch choices whenever their prefixes
agree. The probability of any mismatch is at most $Da$, so the output
total-variation distance from the uniform table law is at most $Da$.
This averages failures at each adaptive history, without conditioning
on a path-dependent event that all count calls succeed.

Set $D_+=\max\{1,D\}$ and choose
\[
 \eta=\gamma=\frac{\tau}{8D_+},\qquad
 t=\left\lceil\log_2\frac{8D_+}{\tau}\right\rceil.
\]
Since $\eta<1/2$, the preceding bound gives $Da\le\tau/2\le\tau$.
There are at most $2D$ count calls and $2D+1$ feasibility tests. Their
input lengths are polynomial in $L$, including the binary encoding of
$\tau$, while $\eta^{-1}=8D_+/\tau$ and
$\log(\gamma^{-1})=\log(8D_+/\tau)$. Theorem~\ref{thm:main} bounds both
the work and the output lengths of every count call, including failed
ones. Thus exact ratio arithmetic and the $t$-bit threshold draw also
have polynomial cost on every execution; the ceiling defining $t$ is
computable by integer shifts and comparisons. The inverse bijection
returns a feasible flow and preserves the total-variation distance.
This proves the sampling and runtime assertions.
\end{proof}

The sampling conclusion is the stated total-variation approximation
with polynomial dependence on $\tau^{-1}$. It asserts neither exact
sampling nor a polynomial-in-$\log(\tau^{-1})$ error dependence.
